\section*{Bab \ref{ch:g12:polymers} --- Polimer}

\begin{solution}{exo:g12:polymers:1}
Propena, \ce{CH2=CH-CH3}.
\end{solution}

\begin{solution}{exo:g12:polymers:2}
\ce{-CF2-CF2-}; \omterm{def:g12:polymers:addition-polymer}{polimer adisi}: \omterm{def:g12:polymers:polymer}{monomernya} memiliki ikatan \ce{C=C} dan
tidak ada produk lain yang terbentuk.
\end{solution}

\begin{solution}{exo:g12:polymers:3}
\omterm{def:g12:curly-arrows:categories}{Adisi}: polietena, PVC, polistirena. Kondensasi: PET (sebuah poliester),
nilon (sebuah poliamida).
\end{solution}

\begin{solution}{exo:g12:polymers:4}
Selulosa, pati, karet dari lateks, protein. Nilon dan PVC adalah \omterm{def:g12:polymers:polymer}{polimer}
sintetis.
\end{solution}

\begin{solution}{exo:g12:polymers:5}
Setiap sambungan memakai satu gugus dari setiap \omterm{def:g12:polymers:polymer}{monomer}; dengan dua
gugus, sebuah \omterm{def:g12:polymers:polymer}{monomer} dapat berikatan di kedua sisinya dan rantainya
terus tumbuh. \omterm{def:g12:polymers:polymer}{Monomer} dengan satu gugus akan menghentikan rantai.
\end{solution}

\begin{solution}{exo:g12:polymers:6}
$M(\ce{C2H3Cl}) = 24.0 + 3.0 + 35.5 = \qty{62.5}{g/mol}$;
$n = 125000/62.5 = 2000$.
\end{solution}

\begin{solution}{exo:g12:polymers:7}
$M(\ce{C8H8}) = \qty{104.0}{g/mol}$; $M = 2500 \times 104.0 =
\qty{2.60e5}{g/mol}$.
\end{solution}

\begin{solution}{exo:g12:polymers:8}
Gugus asam sebuah \omterm{def:g7:atoms-and-molecules:molecule}{molekul} membentuk \omterm{def:g11:functional-groups:carboxylic-acid}{ester} dengan gugus \omterm{def:g11:functional-groups:alcohol}{alkohol} \omterm{def:g7:atoms-and-molecules:molecule}{molekul}
berikutnya, yang masih memiliki gugus asam bebas: rantainya tumbuh. \omterm{def:g12:polymers:repeat-unit}{Unit
ulangnya} \ce{-O-CH(CH3)-CO-}; air dilepaskan pada setiap sambungan.
\end{solution}

\begin{solution}{exo:g12:polymers:9}
Lembaran kantong dibuat dari rantai bercabang, yang tidak dapat tersusun
rapat: bahannya amorf, lunak, dan bening. Botol susu dibuat dari rantai
yang hampir linear, yang tersusun rapat dalam daerah-daerah kristalin:
lebih rapat, lebih kaku, dan daerah kristalin itu menghamburkan cahaya.
\end{solution}

\begin{solution}{exo:g12:polymers:10}
Rantai-rantainya dihubungkan oleh banyak ikatan silang kovalen.
Pemanasan tidak dapat membuat rantai-rantainya bergeser; pemanasan
memutus ikatan dan merusak bahan itu alih-alih melelehkannya.
\end{solution}

\begin{solution}{exo:g12:polymers:11}
$M(\ce{C12H22N2O2}) = 144.0 + 22.0 + 28.0 + 32.0 = \qty{226.0}{g/mol}$;
$n = 22600/226.0 = 100$.
\end{solution}

\begin{solution}{exo:g12:polymers:12}
\[
  \ce{H2N(CH2)6NH2 + HOOC(CH2)4COOH -> H2N(CH2)6NHCO(CH2)4COOH + H2O} .
\]
Produknya masih membawa gugus \omterm{def:g11:functional-groups:amine}{amina} bebas di satu ujung dan gugus asam
bebas di ujung lainnya: masing-masing dapat bereaksi dengan \omterm{def:g12:polymers:polymer}{monomer}
lain.
\end{solution}

\begin{solution}{exo:g12:polymers:13}
Dua \omterm{def:g7:atoms-and-molecules:molecule}{molekul} air per \omterm{def:g12:polymers:repeat-unit}{unit ulang}. $n = 1000/226.0 =
\qty{4.42}{mol}$ \omterm{def:g12:polymers:repeat-unit}{unit ulang}, jadi \qty{8.85}{mol} air:
$8.85 \times 18.0 = \qty{159}{g}$.
\end{solution}

\begin{solution}{exo:g12:polymers:14}
Karet alam yang hangat tersusun atas rantai-rantai yang saling bergeser.
Jembatan belerang menghubungkan rantai-rantai itu: rantai dapat meregang
tetapi kembali lagi, dan tidak lagi mengalir. Karena berikatan silang,
ban tidak dapat dilelehkan dan dibentuk ulang.
\end{solution}

\begin{solution}{exo:g12:polymers:15}
$M(\ce{C6H10O5}) = \qty{162.0}{g/mol}$; $n = \num{1.6e6}/162.0 = 9900$
satuan; panjangnya $9900 \times 0.5 = \qty{4900}{nm}$, sekitar
\qty{5}{\micro m}.
\end{solution}

\begin{solution}{pb:g12:polymers:1}
\textbf{1.} Etana-1,2-diol (dua gugus \omterm{def:g11:functional-groups:alcohol}{alkohol}) dan asam
benzena-1,4-dikarboksilat (dua gugus \omterm{def:g11:functional-groups:carboxylic-acid}{asam karboksilat}).

\textbf{2.} Sambungan \omterm{def:g11:functional-groups:carboxylic-acid}{ester}; air dilepaskan.

\textbf{3.} \omterm{def:g12:polymers:addition-polymer}{Polimer kondensasi}.

\textbf{4.} $M(\ce{C2H6O2}) = \qty{62.0}{g/mol}$; $M(\ce{C8H6O4}) =
\qty{166.0}{g/mol}$.

\textbf{5.} $62.0 + 166.0 - 2 \times 18.0 = \qty{192.0}{g/mol}$; dan
$10 \times 12.0 + 8 \times 1.0 + 4 \times 16.0 = 192.0$.

\textbf{6.} $n = 38400/192.0 = 200$.

\textbf{7.} Dua sambungan \omterm{def:g11:functional-groups:carboxylic-acid}{ester} per \omterm{def:g12:polymers:repeat-unit}{unit ulang}: sekitar 400.

\textbf{8.} Rantai-rantainya sebagian besar linear, dan sebagian dari
rantai itu berjajar dalam daerah-daerah yang teratur.

\textbf{9.} PET adalah \omterm{def:g8:plastics:thermoplastic}{termoplastik}: rantainya tidak berikatan silang
dan saling bergeser jika panas.

\textbf{10.} $300/0.80 = \qty{375}{g}$.

\textbf{11.} $375/25 = 15$ botol.

\textbf{12.} Tutup, label, kotoran, dan sisa potongan disingkirkan atau
hilang selama pemilahan, pencucian, dan pemintalan.

\textbf{13.} Mencegah limbah (prinsip 1): benda bekas menjadi \omterm{def:g5:raw-materials-and-recycling:raw-material}{bahan
baku}.

\textbf{14.} \ce{C10H8O4 + 2H2O -> C2H6O2 + C8H6O4}.

\textbf{15.} $1000/192.0 = \qty{5.21}{mol}$.

\textbf{16.} Diol: $5.21 \times 62.0 = \qty{323}{g}$; asam:
$5.21 \times 166.0 = \qty{865}{g}$.

\textbf{17.} $2 \times 5.21 \times 18.0 = \qty{188}{g}$ air;
$1000 + 188 = 323 + 865 = \qty{1188}{g}$: massa kekal.

\textbf{18.} \omterm{def:g12:polymers:polymer}{Monomer-monomernya} dapat dimurnikan dan menghasilkan PET
baru yang sebaik PET baru, sedangkan pelelehan memperpendek rantainya
sedikit demi sedikit setiap kali.

\textbf{19.} Satu jaket bulu sintetis memerlukan 15 botol \qty{25}{g}.
\end{solution}
